%%
%% This is file `example/ch-intro.tex',
%% adapted for the current thesis content.
%%

\chapter{绪论}

\section{研究背景与意义}

从功能性磁共振成像（functional Magnetic Resonance Imaging, fMRI）信号中重建人类的视觉体验，是计算神经科学与脑机接口领域的重要研究方向。该问题旨在揭示大脑皮层中所编码的视觉信息，并建立从脑活动到感知内容的直接映射关系。一旦实现高质量的视觉重建，不仅能够为理解人类视觉认知机制提供新的研究工具，还将为基于非侵入式脑信号的神经接口系统提供关键技术支撑。

近年来，随着深度学习和生成模型的快速发展，基于神经解码的视觉重建技术取得了显著进展。特别是预训练视觉表示模型（如 CLIP）与扩散生成模型的引入，使得重建图像在语义一致性与视觉逼真度方面均有明显提升。然而，现有大多数方法仍然局限于单被试（within-subject）设置，即需要针对每个被试分别收集大量配对的 fMRI-图像数据进行训练或微调。这种高度依赖个体数据的模式不仅成本高昂，而且严重限制了模型的可扩展性和实际应用能力。

为提升模型的泛化能力，研究者逐渐将关注点转向跨被试（cross-subject）视觉重建。其中，更具挑战性的情形是零样本跨被试（zero-shot cross-subject）重建，即模型仅在一组源被试上训练，然后直接应用于完全未见过的目标被试，且不进行任何微调或校准。这一设定能够更真实地反映模型的泛化能力，但同时也带来了显著困难。

其核心挑战在于：fMRI 信号中同时包含两类纠缠的信息，一类是由外界视觉刺激驱动的语义信息，另一类则是由个体差异（如脑结构、功能组织及认知策略）以及扫描噪声、设备误差等引入的被试特异性变化。当这两类信息未被有效分离时，模型往往会过度依赖被试特征，从而削弱对语义信息的提取能力。此外，配对 fMRI-图像数据本身极为稀缺，使得在数据受限条件下学习被试不变的表示更加困难。

因此，研究如何在有限数据条件下，从多被试 fMRI 信号中提取具有良好跨个体泛化能力的语义表征，对于推动视觉脑解码技术的发展具有重要的科学意义和应用价值。

\section{国内外研究现状}

\subsection{单被试视觉重建方法}

早期的 fMRI 视觉重建研究主要集中在单被试场景，通常基于编码—解码框架或贝叶斯方法，实现从脑活动到视觉内容的反演。这类方法验证了从视觉皮层活动中恢复感知信息的可行性。随着深度学习的发展，研究者逐渐采用深层神经网络提取层次化视觉特征，以更好地匹配大脑皮层的表征结构，从而显著提升重建质量。

近年来，多模态表示学习和生成模型的引入进一步推动了该方向的发展。通过结合视觉—语言对齐模型（如 CLIP）与扩散模型，可以在高维语义空间中进行解码，并生成具有较高语义一致性的图像。代表性方法如 MindEye、NeuroPictor 等，已在多个标准评测指标上取得较好性能。然而，这类方法高度依赖被试特定数据，难以推广至新个体。

\subsection{跨被试视觉解码方法}

为减少对个体数据的依赖，跨被试脑解码方法尝试学习在不同被试之间共享的神经表征。这类方法通常基于潜空间对齐、功能对齐或被试不变建模等技术，以缓解个体差异带来的影响。

近年来，一些方法结合大规模视觉先验与 Transformer 结构，在多被试数据上学习通用映射关系，并通过少量校准数据对新被试进行适配。这类方法在降低数据需求的同时，能够在一定程度上保持性能。然而，由于仍依赖少量目标被试数据，其泛化能力在严格无监督条件下仍然受限。

\subsection{零样本跨被试重建方法}

零样本跨被试重建进一步提高了问题难度，要求模型在推理阶段完全不依赖目标被试数据。这一设定能够更纯粹地衡量模型学习被试不变语义表示的能力。已有研究表明，通过显式对齐和特征解耦，可以在一定程度上提升跨被试泛化性能。

尽管如此，与有监督或少样本设置相比，当前零样本方法仍存在明显性能差距。在被试差异与刺激分布同时变化的情况下，模型性能下降尤为显著。已有研究指出，仅依赖数据规模扩展或强生成先验并不足以解决该问题，仍需在模型层面引入针对被试不变性、语义一致性及表征解耦的显式建模机制。

\subsection{现有方法的不足}

综上所述，现有研究主要存在以下问题：

\begin{itemize}
  \item 依赖被试特定数据，泛化能力有限；
  \item 缺乏对被试差异的显式建模；
  \item 语义信息与个体差异存在严重纠缠；
  \item 数据稀缺限制了跨被试表示学习的效果。
\end{itemize}

因此，如何在数据受限条件下构建具有良好跨被试泛化能力的视觉解码模型，仍是亟待解决的关键问题。

\section{本文主要工作概述}

针对零样本跨被试 fMRI 视觉重建中存在的被试差异显著、语义信息与个体特征相互纠缠以及训练数据稀缺等问题，本文从跨被试表示学习的角度出发，提出了一种统一的端到端重建框架 MINDMELD。该方法通过引入神经响应合成与语义解耦机制，在扩展训练数据分布的同时学习具有跨被试泛化能力的语义表示。

首先，为缓解配对 fMRI-图像数据稀缺所带来的监督不足问题，本文构建了一种被试感知的神经响应合成模型。具体而言，该模型以视觉刺激为输入，利用预训练视觉特征提取网络提取多层语义信息，并结合被试的群体感受野（population receptive field, pRF）特征生成被试相关的调制参数，对共享的皮层顶点表示进行条件调制，从而预测不同被试在相同刺激下的皮层响应。该设计在保持视觉语义一致性的同时，引入了可控的被试差异，从而有效扩展了神经响应的分布范围。

在此基础上，为解决神经响应中语义信息与被试特异性信息难以分离的问题，本文设计了一种被试不变语义解码模型。该模型首先将顶点级神经响应编码为结构化的潜在表示，并进一步分解为语义表示与被试表示两个分支。其中，语义分支通过与图像语义空间对齐来捕获共享的视觉信息，而被试分支用于建模个体差异。为了防止语义表示中残留被试信息，本文在语义分支上引入梯度反转层（Gradient Reversal Layer, GRL）及对抗分类器，通过对抗学习机制显式抑制被试身份信息，从而提升语义表示的跨被试泛化能力。

进一步地，为增强不同被试之间的语义对齐能力，本文提出了一种跨被试语义一致性约束机制。具体而言，对于多个被试在相同视觉刺激下的神经响应，首先在语义空间中进行聚合以形成共享语义表示，然后结合各自的被试特征进行重建，使其能够恢复对应的神经潜表示。通过该跨重建过程，模型被迫学习在不同被试之间保持一致且具有解释能力的语义表示，从而显式强化跨被试语义一致性。

综上所述，本文的主要工作可概括如下：

\begin{enumerate}
  \item 提出一种结合视觉特征与被试 pRF 信息的神经响应合成方法，用于生成多被试条件下的皮层响应，从而扩展训练数据分布；
  \item 构建基于特征解耦的被试不变语义解码模型，并通过对抗学习抑制语义表示中的被试信息，实现跨被试语义建模；
  \item 设计跨被试语义一致性约束，通过跨重建机制强化不同被试间的语义对齐能力，提升零样本条件下的泛化性能。
\end{enumerate}

\chapterbib
